\section{Verification and measured performance}
\label{sec:experiments}

The algebraic separation theorems and the observed performance lead to
different conclusions. Corridor transfer removes provably unnecessary
work from the inherited forward transfer, and its sparse scalar variant
removes a quadratic representation cost on the explicit separation
family. On the ordinary diagram scans measured here, however, the
production sparse reducer remains substantially faster. This section
reports both conclusions, with the controls needed to interpret them.

\subsection{What the verification checks}

The algebraic tests compare several levels of information. Before and
after ordinary cancellation, the full
\((\text{matching},h,q)\) survivor profile is checked against scalar
homology, together with quantum homogeneity and \(d^2=0\). On allocated
prefixes, the new transfer is compared entrywise with the frozen
one-sided implementation in the same scalar-contraction basis, in
forward, reverse, and automatic modes, with pruning both enabled and
disabled. Selected cases also check all eight full contraction
identities in the cobordism algebra. These checks include parity
cancellation: a nonempty corridor can have zero path sum over \(\F\).
Boolean reachability is never interpreted as a nonzero coefficient.

An independent set-valued coefficient engine with a different pivot
order checks the final homological-degree ranks on seeded diagram
families. Sparse scalar-component tests convert every tuple column
back to a binary mask and compare inclusion, projection, and homotopy
entrywise with the frozen scalar construction. Further tests exercise
dead object slots, partially completed ordinary cancellation, tail
contraction, the component coefficient engine, and resource failure.
The complete scanner suite passes \(268\) tests: \(245\) baseline
tests, eight graded-support tests, ten corridor tests, and five
scalar-component tests. Its recorded run took \(30.764\) seconds;
the package retains the log and final verification manifest.
These finite checks support the implementation; they do not replace
the general proofs.

The geometric kernel has a separate expanded-sheet oracle. It
constructs the lifted graph, finds its connected components by graph
search, propagates orientation parity along lifted edges, and
enumerates boundary permutations. The oracle does not reuse the
classifier's gcd calculation, boundary-word evaluator, or canonical
schema constructor. The supplied campaign comprises \(13{,}846\)
topology comparisons, including exhaustive low-rank assignments and
\(2{,}000\) seeded larger schemas. Additional checks cover fibre
relabeling, explicit generic equivariance, exhaustive bijections for
small exceptional covers, malformed inputs, and small moduli.

The large-integer test uses the reflection action over a M\"obius
band with \(W=2^{10000}\) sheets. It returns three family records,
two cover-isomorphism types, and \(2^{9999}+1\) connected components
without expanding them. This verifies the intended compressed output
contract; it is not a benchmark of normal-surface extraction or
unknot recognition.

\subsection{Timing protocol and baselines}
\label{sec:timing-protocol}

The principal measurements are recorded in
\path{fast/results/corridor_benchmark_20261008.json}, generated by
\path{fast/benchmark_corridor.py}. The software environment recorded
in the file is Python \(3.12.14\) on Linux \(6.18.44\), x86-64, with
glibc \(2.39\). Each case has seven paired rounds, with the order of
the arms shuffled by a fixed seed. Every implementation is warmed
before measurement. A warm pilot selects a fixed repetition count
for each arm, targeting a \(20\)-millisecond batch, capped at
\(32\) repetitions. The two identical sparse-control arms use the
same repetition count.

For actual diagrams, a fresh scanner is constructed on every
repetition and the complete raw scan is timed. All arms receive the
same precomputed crossing order and the same scanner-level
shape-cache setting. Diagram construction, common order selection,
and explicit verification are excluded. No braid, Seifert,
Alexander, Jones, Reidemeister, or factorization front end is being
credited with a speedup in this experiment. For synthetic complexes,
construction is excluded and the complete reduction of a fresh
allocated complex is timed. Scalar contraction, graph construction,
pruning, direction selection, and output installation are included.

Before each measured batch the garbage collector is collected; it
remains enabled during timing. Every scanner has a cooperative
\(30\)-second deadline. Other team agents paused CPU-intensive work
during the final coordinated runs. Earlier exploratory timing files
are retained as provenance but are not the primary dataset.

The nine arms are the production sparse scanner, an identical sparse
control, quantum bookkeeping with the ordinary reducer, the frozen
one-sided transfer, forward corridor evaluation without pruning,
forward evaluation with pruning, reverse evaluation with pruning,
automatic direction selection with pruning, and adaptive corridor
reduction. The automatic full-transfer arm uses the original global
scalar construction and packed endpoint reach sets. The later
sparse-component/Boolean variant is measured separately.

All times in the tables are medians of per-run batch times. Ratios
are computed within each paired round and then summarized by their
median; consequently a reported ratio need not equal the ratio of
the two displayed median times. An \(A/A\) ratio compares the two
identical sparse implementations. Across the principal cases these
ratios range from \(0.755\) to \(1.399\), despite the controls.
Absolute timings and small differences therefore deserve
considerable caution. We report large deterministic count
reductions separately from elapsed-time ratios, and do not promote
small apparent gains into general performance claims.

\subsection{A negative result for the stronger support certificate}

The full grading allows the zero-differential certificate of
Proposition~\ref{prop:corridor-typed} even when consecutive
homological degrees survive. The audit
\path{fast/audit_graded_support.py} evaluates this certificate and
the earlier degree-gap certificate on exactly the same unreduced
prefixes, then checks the predicted profile against ordinary
cancellation.

Its \(72\) actual diagrams consist of twelve fixed fixtures or
families and sixty seeded random braid closures. Across \(639\)
stages, both certificates succeed on exactly \(196\) stages; the
full graded certificate supplies \emph{zero additional successes}.
The ordinary reductions perform \(3{,}470\) Schur-update pairs.
The complete records are in
\path{fast/results/graded_support_20261008.json}.

This does not weaken the support theorem. It shows that a strictly
stronger sufficient condition need not improve the supplied
execution family. It also explains why the implementation effort
proceeds to complete transfer, which retains nonzero radical
differentials, instead of assuming that a better all-zero test will
usually suffice.

\subsection{Actual raw diagram scans}

Table~\ref{tab:actual-corridor} gives the main actual-diagram result.
The first case is the closure of the \(201\)-letter positive
two-braid. The next two are the closures of
\((\sigma_1\sigma_2)^{10}\) and
\((\sigma_1\sigma_2^{-1})^4\). The remaining cases use the supplied
Morton, Conway, and hard eight-crossing fixtures. The exact PD codes
and crossing orders are recorded in the JSON.

\begin{table}[htbp]
\centering
\small
\setlength{\tabcolsep}{4pt}
\begin{tabular}{@{}lrrrrrr@{}}
\toprule
Case & Sparse ms & Full ms & Adaptive ms &
Sparse/full & Sparse/adapt. & \(A/A\)\\
\midrule
Two-braid, \(201\) & 124.044 & 1324.193 & 192.852 & 0.094 & 0.679 & 0.913\\
Torus \(3,10\) & 8.936 & 29.650 & 9.692 & 0.311 & 0.954 & 1.268\\
Alternating three-braid & 5.059 & 14.435 & 4.705 & 0.367 & 1.141 & 1.009\\
Morton & 2.075 & 4.801 & 2.351 & 0.412 & 1.026 & 0.947\\
Conway & 5.828 & 17.704 & 6.937 & 0.319 & 0.864 & 0.878\\
Hard unknot, \(8\) & 1.195 & 2.796 & 1.335 & 0.475 & 0.934 & 0.938\\
\bottomrule
\end{tabular}
\caption{Complete raw scans. \emph{Full} means automatic pruned corridor
transfer at every reduction stage. A ratio greater than one would
favor the corresponding corridor arm.}
\label{tab:actual-corridor}
\end{table}

Full corridor reduction is approximately \(2.1\)--\(10.7\) times
slower than the production sparse class on these cases. It can
remove a substantial fraction of the graph while still losing
overall, because binary contraction and graph setup must be paid
at every stage. This is a measured total-cost failure, not a
failure of coefficient correctness: all final ranks by homological
degree agree.

The adaptive arm performs no corridor switches on any of these six
cases. Its sparse-work allowance is never exceeded. Thus the modest
variations in its timing cannot be attributed to successful use of
corridor transfer. The data support retaining ordinary sparse
cancellation as the default; they supply no reliable general
actual-diagram speedup from either new strategy.

\subsection{Homogeneous algebraic families with nonzero transfer}

The synthetic tests are valid two-term complexes in the same
graded category. Since they occupy only homological degrees zero
and one, their square-zero condition is automatic; quantum
homogeneity is checked explicitly. They are not asserted to occur
as prefixes of any supplied classical knot diagram.

The seeded dense family uses one matching with three arcs.
Surviving sources have quantum shift \(0\), surviving targets have
shift \(6\), and two scalar contractible layers have shifts \(2\)
and \(4\). Radical matrices between successive layers are
multiplied by \(x,y,z\), respectively. Each scalar layer has
sixteen contractible pairs. Source-heavy and target-heavy cases
use multiplicities \(64:1\) and \(1:64\); the balanced case uses
\(16:16\). The dense dead-branch case adds sixty-four contractible
pairs at each intermediate shift on branches that cannot reach a
surviving target.

The strict shared-suffix cases use the sparse family of
Theorem~\ref{thm:corridor-separation}, with \(R=M\) odd. The
dead-leaf case uses Proposition~\ref{prop:corridor-dead} with
\(255\) irrelevant contractible leaves. Every displayed strict
case has nonzero transferred differential. Thus the reductions
cannot be explained by an all-zero shortcut.

\begin{table}[htbp]
\centering
\small
\setlength{\tabcolsep}{5pt}
\begin{tabular}{@{}lrrrr@{}}
\toprule
Family & Prior attempts & Corridor attempts &
Prior/corridor & Sparse/corridor\\
\midrule
Dense source-heavy & 3,940 & 307 & 3.631 & 0.506\\
Dense target-heavy & 317 & 317 & 0.365 & 0.243\\
Dense balanced & 2,142 & 2,142 & 1.070 & 0.397\\
Dense dead branches & 72,204 & 342 & 9.690 & 1.346\\
Shared suffix, \(15\) & 465 & 45 & 1.975 & 0.157\\
Shared suffix, \(63\) & 8,001 & 189 & 7.091 & 0.105\\
Shared suffix, \(255\) & 130,305 & 765 & 28.409 & 0.098\\
Dead leaves, \(255\) & 256 & 1 & 0.498 & 0.120\\
\bottomrule
\end{tabular}
\caption{Radical-edge product attempts and paired total-reduction
timing ratios. The prior is the frozen one-sided transfer; the
corridor arm uses automatic direction and pruning. Synthetic
construction is excluded, but scalar and graph setup are included.}
\label{tab:kernel-corridor}
\end{table}

The operation counts in Table~\ref{tab:kernel-corridor} compare the
same event: propagation across a radical edge, requiring a typed
coefficient product. They do not compare the predecessor's raw
algebra-call counter directly with the new implementation's
composition counter. The latter excludes typed identity shortcuts
and may use cached products. Scalar propagations are distinguished
from radical attempts; vertex-loop and setup costs remain included
in the elapsed times.

At \(R=M=255\), the exact formulas give
\[
 R(1+2M)=130{,}305,\qquad R+2M=765.
\]
The radical-edge reduction factor is \(170\frac13\), while the
measured total reduction is about \(28.4\) times faster than the
inherited forward sweep. Ordinary Markowitz cancellation is still
about ten times faster than the full corridor method on this
family. Its pivot order already exploits the shared structure
effectively. The theorem proves an unbounded improvement over a
specified inherited transfer, not superiority over every exact
reduction strategy.

The dead-leaf example makes the setup issue especially clear:
radical attempts drop from \(256\) to \(1\), yet total reduction
becomes slower than the prior. Conversely, the dense dead-branch
case has a modest apparent \(1.35\)-fold gain over sparse reduction.
Given the identical-control drift, we do not treat that single
ratio as dependable practical superiority.

All full-transfer arms produce exactly the same minimal
differential in the frozen scalar basis; the benchmark records its
SHA-256. Ordinary and adaptive sparse reductions may select a
different homogeneous basis. They are checked independently by
survivor multiplicities and by the binary rank of the coefficient
matrix of the single possible output monomial, \(xyz\) or \(x\).
For this planted family those quantities are preserved by the
allowed homogeneous changes of basis.

\subsection{The sparse scalar refinement: a separate negative timing test}

After the principal experiment was frozen, a targeted follow-up combined
sparse scalar components, Boolean reachability, and direction selection
by port count. It compared this variant with the frozen corridor's
forward and automatic modes, the port-count policy with
global scalar setup, and two identical production sparse arms.
This keeps the effect of the new setup representation separate from
the original direction/pruning result. Its data are in
\path{fast/results/compressed_corridor_20261008.json}.

\begin{table}[htbp]
\centering
\small
\setlength{\tabcolsep}{5pt}
\begin{tabular}{@{}lrrrrr@{}}
\toprule
Case & Frozen auto ms & New sparse ms &
Auto/new & Production/new & \(A/A\)\\
\midrule
Two-braid, \(201\) & 1269.243 & 1448.700 & 0.803 & 0.091 & 1.087\\
Conway & 20.762 & 26.421 & 0.804 & 0.292 & 1.110\\
Hard unknot, \(8\) & 2.746 & 2.866 & 0.934 & 0.441 & 0.795\\
Shared suffix, \(15\) & 0.562 & 0.423 & 1.250 & 0.174 & 1.191\\
Shared suffix, \(63\) & 1.672 & 1.422 & 1.202 & 0.132 & 1.056\\
Shared suffix, \(255\) & 7.847 & 8.178 & 0.871 & 0.101 & 1.011\\
Dead leaves, \(255\) & 2.663 & 2.392 & 1.113 & 0.154 & 0.997\\
\bottomrule
\end{tabular}
\caption{The targeted sparse-setup follow-up. \emph{New sparse} means
sparse scalar components, Boolean pruning, and port-count direction
selection, not the production sparse cancellation algorithm.}
\label{tab:compressed-setup}
\end{table}

There is no stable practical timing gain at these tested sizes. The
refinement loses on all three actual examples and on the largest
shared-suffix instance; the small apparent gains elsewhere are
comparable to control variation. Its justification is the exact
representation and complete-stage bound, not a claim that Python
tuples and extra component bookkeeping are always faster.

\subsection{Exact representation counts}

The additional audit
\path{fast/audit_scalar_representation.py} compares all scalar
matrices entrywise and records representation sizes in
\path{fast/results/scalar_representation_20261008.json}. In the
shared-suffix constructor's fixed object order, the sum of the
integer bit lengths of all inherited \(i,p,h\) columns is
\[
 R^2+4R+5+MR+\frac{M^2+9M}{2}.
\]
At \(R=M\), this is \((5M^2+17M+10)/2\). The sparse representation
has exactly \(2R+M+3\) nonzero indices, each requiring only the
binary length of an object or survivor index.

For completeness, the constructor orders the sources first, then
\(L_0,B_0\), then all \(L_i\), then all \(B_i\), and finally
\(t\). The separate sums of bit lengths are
\[
\begin{aligned}
 |i|_{\rm bits}&=R(R+1)/2+R+2M+3,\\
 |p|_{\rm bits}&=R(R+1)/2+R+1,\\
 |h|_{\rm bits}&=R+1+\sum_{i=1}^{M}(R+i+2).
\end{aligned}
\]
Adding proves the displayed polynomial. There are respectively
\(R+1\), \(R+1\), and \(M+1\) nonzero scalar entries, proving the
sparse count. Other object orderings can change the packed-bit
polynomial, while the sparse nonzero count is unchanged for this
direct contraction.

\begin{center}
\small
\begin{tabular}{@{}rrr@{}}
\toprule
\(R=M\) & Packed integer bit length & Sparse indices\\
\midrule
1 & 16 & 6\\
3 & 53 & 12\\
15 & 695 & 48\\
63 & 10,463 & 192\\
255 & 164,735 & 768\\
\bottomrule
\end{tabular}
\end{center}

For \(M=255\), there are \(768\) original objects and \(512\)
scalar components: \(256\) singletons and \(256\) identity pairs.
No component requires nontrivial binary elimination. This is the
situation in which sparse setup makes the complete-stage separation
possible. The table measures integer payload length and support
counts, not actual Python heap allocation; object and tuple
overheads can dominate at small sizes.

The experiments therefore leave a clear deployment decision.
Preserve the production sparse default, retain the new kernels as
exact opt-in research capabilities, and use the supplied counters
to identify inputs where setup is amortized. A universal
quasi-polynomial bound still requires the structural hypotheses
discussed in Section~\ref{sec:accounting}.
